\scanpage{143}

% L1, introduction
Récapitulons

% L2–6, theorem and hypotheses
\uline{Thm} $(E,q)$ forme quadratique régulière de $\mathrm{rg}\ E=N$, $N=2n\deleted{\unsure{+1}}$ ou $N=2n+1$ ($n\geq1$ donc $N\geq2$). Soient $H,H'\subset[1,n]$, $\underset{\deleted{\unsure{d}}}{d}=\max H$, $d'=\max H'$ (donc $0\leq d,d'\leq n$), considérons

% L7, numbered display
\[
\underline\Delta_H(E,q)\xrightarrow{\unsure{p|\rho}_{H',H}}\underline\Delta_{H'}(E,q)\tag*{(1)}
\]

% L8, clause a
a) Ce morphisme est projectif, \deleted{et} lisse, \deleted{\unsure{à}} \inserted{surjectif}.

% L9–10, consequence
\deleted{b)} En particulier (prenant $H'=\varnothing$) les $\underline\Delta_H(E,q)\to S$ sont projectifs, lisses, surjectifs.

% L11–14, clause b
b) Les fibres géométriques de ce morphisme sont \deleted{\unsure{0}} connexes, sauf dans le cas $N$ pair, $d'<d=n$, auquel cas les fibres géométriques \deleted{\unsure{sont|ont}}\unsure{ont|sont} deux comp conn.

% L15, clause c
c) Dans le \deleted{dernier cas} \inserted{cas $N=2n$}, considérons

% L16–17, display
\[
\begin{aligned}
\underline{\mathrm{Arf}}(E,q)={}&\text{schém. des comp. conn}\\
&\text{de }\underline\Delta_{\{n\}}(E,q)\ \text{sur }S
\end{aligned}
\]

% L18–19, lower text
qui est un rev. étale de \deleted{degré} 2 de $S$. Considérons \inserted{(si $H\supset\{n\}$)} le \deleted{\unsure{fonct.|facteur}} morphisme composé

% L20, lower display
\[
\underline\Delta_H\longrightarrow\underline\Delta_{\{n\}}\longrightarrow\underline{\mathrm{Arf}}
\]

% L21, bottom text
\deleted{C'est lui qui \unsure{est ceci|est ici|ratisfait}} Ce morphisme est

% bottom left, pencil
% archive sheet number: 142
